\section*{第 \ref{ch:b2:plant-meristems} 章——植物营养发育：分生组织与生长}

\begin{solution}{exo:b2:plant-meristems:1}
\omterm{def:b2:plant-meristems:meristem}{分生组织}是持续分裂的未\omterm{def:b2:cell-differentiation:differentiation}{分化}细胞群，
植物的组织都由它衍生。\omterm{def:b2:plant-meristems:meristem}{茎尖分生组织}：
形成叶、茎、花；\omterm{def:b2:plant-meristems:meristem}{根尖分生组织}：形成根；\omterm{def:b2:plant-meristems:wood}{维管
形成层}：形成次生木质部（\omterm{def:b2:plant-meristems:wood}{木材}）和韧皮部；木栓\omterm{def:b2:plant-meristems:meristem}{形成层}：
形成\omterm{def:b2:plant-meristems:wood}{树皮}的\omterm{def:b2:plant-meristems:wood}{周皮}（木栓）。
\end{solution}

\begin{solution}{exo:b2:plant-meristems:2}
\omterm{prop:b2:plant-meristems:ram}{根冠}（保护、润滑、感受重力，并不断脱落细胞）；
分生区（含\omterm{prop:b2:plant-meristems:ram}{静止中心}，负责分裂）；\omterm{prop:b2:plant-meristems:ram}{伸长区}
（细胞伸长十至二十倍）；\omterm{def:b2:cell-differentiation:differentiation}{分化}区（\omterm{prop:b2:plant-meristems:ram}{根毛}、成熟的木质部
和韧皮部）。侧根起源于成熟根的中柱鞘，
与木质部极相对。
\end{solution}

\begin{solution}{exo:b2:plant-meristems:3}
将豆株去顶，侧芽随之长出。另取一组，
在切口处放置含\omterm{prop:b2:plant-meristems:apical}{生长素}的琼脂块，侧芽保持\omterm{def:b2:angiosperm-reproduction:seed}{休眠}；
对照组放置普通琼脂块，侧芽则生长。
由此可知，来自顶端的\omterm{prop:b2:plant-meristems:apical}{生长素}抑制侧芽，
即\omterm{prop:b2:plant-meristems:apical}{顶端优势}是由激素调节的。
\end{solution}

\begin{solution}{exo:b2:plant-meristems:4}
六十岁；边材是最外层的十五个\omterm{def:b2:plant-meristems:wood}{年轮}；活的部分有
\omterm{def:b2:plant-meristems:meristem}{形成层}、韧皮部、射线细胞和边材的薄壁细胞，
以及木栓\omterm{def:b2:plant-meristems:meristem}{形成层}，后者是围绕死亡心材的一层薄壳。
\end{solution}

\begin{solution}{exo:b2:plant-meristems:5}
$0.1\times(P - 0.3)$：\qty{0.01}{h^{-1}}，倍增时间 $\ln 2/0.01 =
\qty{69}{h}$), \qty{0.03}{h^{-1}} (\qty{23}{h}), \qty{0.05}{h^{-1}}
（\qty{14}{h}）。在 \qty{0.25}{MPa} 时，膨压低于屈服
阈值，生长停止，早于任何萎蔫出现。
\end{solution}

\begin{solution}{exo:b2:plant-meristems:6}
叶 1--8 的位置角依次为 0、137.5、275、52.5、190、327.5、105、\qty{242.5}{\degree}。
排序后为 0、52.5、105、137.5、190、242.5、275、327.5——最小间隔
为 \qty{32.5}{\degree}。每片叶都落在前面各叶所形成的一个间隙中，
从而八片叶彼此遮挡尽可能少，
树冠便均匀地铺满受光圆盘。
\end{solution}

\begin{solution}{exo:b2:plant-meristems:7}
$2\pi r h\Delta r$：在 \qty{5}{cm} 处，$2\pi\times 0.05\times 12\times
0.003 = \qty{0.011}{m^3}$，\qty{5.7}{kg}，其中碳 \qty{2.8}{kg}；在
\qty{40}{cm} 处，\qty{0.090}{m^3}，\qty{45}{kg}，其中碳 \qty{23}{kg}
——同样的环宽下，碳量为前者的八倍。
\end{solution}

\begin{solution}{exo:b2:plant-meristems:8}
\emph{clv3}：缺少“刹车”——\omterm{def:b2:plant-meristems:meristem}{分生组织}极度膨大，多出的叶片和花
器官，茎发生带化。\emph{wus}：缺少“加速器”——\omterm{def:b2:plant-meristems:meristem}{分生组织}
仅产生几片叶便耗尽，并反复重建。双\omterm{def:b2:genome-diversification:mutation}{突变}体的
\omterm{def:b2:meiosis-heredity:vocabulary}{表型}与 \emph{wus} 相同，因此 WUS 位于 CLV3 的下游：
CLV3 的作用是抑制 WUS，若无 WUS，
便无可抑制。
\end{solution}

\begin{solution}{exo:b2:plant-meristems:9}
运输：\qty{20}{h}。扩散：$(0.2)^{2}/(2\times 10^{-9}) =
\qty{2e7}{s}$，约 230 天。\omterm{prop:b2:plant-meristems:apical}{极性运输}使激素信号
可在一天之内沿整株植物的长度发挥作用，并赋予
其方向。
\end{solution}

\begin{solution}{exo:b2:plant-meristems:10}
成熟组织每天 \qty{3}{cm}，细胞长 \qty{200}{\micro m}：
每条细胞列每天 \num{150} 个细胞。故一百个分裂细胞
须每天各分裂 1.5 次，周期约 \qty{16}{h}。
\end{solution}

\begin{solution}{exo:b2:plant-meristems:11}
赤霉素使节间伸长；不能响应赤霉素的植株
节间短、茎也短。重施氮肥使高秆小麦
长得更高、穗更重，直至倒伏
并腐烂；矮而挺的茎秆能直立，并把多余的
碳更多地投入籽粒而非秸秆，因而收获指数上升。在
野外，矮株被邻近植株压过、遮荫，
在竞争中落败。
\end{solution}

\begin{solution}{exo:b2:plant-meristems:12}
\omterm{def:b2:plant-meristems:meristem}{分生组织}使植物能在条件适宜的时间与地点增添器官
——叶向光、根向水——并舍弃
那些不划算的器官；每个腋芽都是备用生长点，
因此啃食、火烧和修剪之后，植株可凭借再生长存活；并且
由于\omterm{def:b2:plant-meristems:meristem}{分生组织}始终保持胚性，\omterm{def:b2:asexual-reproduction:clone}{基株}没有
固定寿命。代价是：植物不能移动，也不能重组已建成的
结构——器官一旦着生便固定原处，而植体以累积方式
建成，须带着其已死的部分（心材）。
\end{solution}

\begin{solution}{pb:b2:plant-meristems:1}
\textbf{1.} $150/3 = 50$，每个顶端有 50 片叶；全树 \num{10000} 片。
\textbf{2.} 0、137.5、275、52.5 和 \qty{190}{\degree}。
\textbf{3.} $8\times 137.5 = 1100 = 3\times 360 + 20$; $13\times 137.5
= 1787.5 = 5\times 360 - 12.5$; $21\times 137.5 = 2887.5 = 8\times 360
+ 7.5$：第 22 片叶是第一个与首片
相距在 \qty{10}{\degree} 以内的叶。
\textbf{4.} 每片新叶都落在最宽的间隙中，因此叶片彼此
遮挡最少；决定间距的是\omterm{prop:b2:plant-meristems:apical}{生长素}，其\omterm{prop:b2:plant-meristems:apical}{极性运输}（依靠 PIN
载体蛋白）抽离每个原基周围的外界，
从而决定间距。
\textbf{5.} $\num{10000}\times 80\,\text{cm}^{2} = \qty{80}{m^2}$：整个
树冠即为本季的新叶。
\textbf{6.} 全部如此；最初的叶片由前一年
贮存在\omterm{def:b2:plant-meristems:wood}{木材}与根中的淀粉构成。
\textbf{7.} 顶端增大为三倍后原基形成加快——每个
\omterm{prop:b2:plant-meristems:sam}{叶序}周期出现两到三个；叶片与叶面积可增至三倍，但枝条
将发生带化，茎呈带状，\omterm{prop:b2:plant-meristems:sam}{叶序}紊乱，
且大部分多出的叶片会自相遮荫。
\textbf{8.} 白天：$0.05\times(0.55 - 0.35) = \qty{0.01}{h^{-1}}$；夜晚：
$0.05\times 0.40 = \qty{0.02}{h^{-1}}$.
\textbf{9.} 白天：$2e^{0.12} = \qty{2.25}{cm}$，增加 \qty{0.25}{cm}；
夜晚：$2.25e^{0.24} = \qty{2.87}{cm}$，增加 \qty{0.61}{cm}。
\textbf{10.} 在 \qty{96}{h} 内平均速率为 \qty{0.015}{h^{-1}}：$2e^{1.44}
= \qty{8.4}{cm}$.
\textbf{11.} 白天：$0.05\times 0.30 = \qty{0.015}{h^{-1}}$；夜晚：
$0.05\times 0.50 = \qty{0.025}{h^{-1}}$.
\textbf{12.} $0.05\times 0.05 = \qty{0.0025}{h^{-1}}$，为
白天正常速率的四分之一。渗透调节提高细胞内溶质含量，
使得在相同水势下，膨压重新升至
阈值之上。
\textbf{13.} 白天蒸腾作用降低叶片的水势，
膨压随之下降，故 $P - Y$ 很小；夜晚膨压
得以恢复。生长逐时受水分限制，而非受糖限制。
\textbf{14.} $2\pi\times 0.10\times 8\times 0.005 = \qty{0.025}{m^3}$;
干重 \qty{12.6}{kg}。
\textbf{15.} 其中碳 \qty{6.3}{kg}。
\textbf{16.} $80\times 4\times 150 = \qty{48}{kg}$，为碳。
\textbf{17.} $6.3/48 = \qty{13}{\%}$。其他库包括叶片
自身、根、枝条及小枝，每个活细胞的呼吸作用
（约占所固定碳的一半）、储备物、花与
\omterm{def:b2:angiosperm-reproduction:seed}{种子}。
\textbf{18.} 半径 \qty{0.20}{m}：$2\pi\times 0.20\times 8\times
0.005 = \qty{0.050}{m^3}$，为今年的两倍——\omterm{def:b2:plant-meristems:meristem}{形成层}所沿
的周长已加倍（树也将更高）。
\textbf{19.} 干旱或寒冷的年份；或虫害导致落叶、霜冻
或火灾。气候因素会使该地区每棵树的同一\omterm{def:b2:plant-meristems:wood}{年轮}
同时变窄，并与气象记录相关；而损伤只影响一棵树
或一片林分，霜冻或火灾则会在\omterm{def:b2:plant-meristems:wood}{年轮}中留下解剖学伤痕。
\textbf{20.} 靠近每处切口的芽在数日内萌发，长成
数个枝条；一月后树冠更丰满、更稠密，枝条更多，
但每根更短。
\textbf{21.} 切口处的\omterm{prop:b2:plant-meristems:apical}{生长素}替代了顶端的信号：芽
保持\omterm{def:b2:angiosperm-reproduction:seed}{休眠}，不再出现分枝。
\textbf{22.} 每次修剪都去除顶端，释放
下方的芽；一年六次，每次都增加一层短枝，
树冠表面便布满了这些短枝。
\textbf{23.} 来自树桩上的\omterm{def:b2:angiosperm-reproduction:seed}{休眠}腋芽，以及
\omterm{def:b2:plant-meristems:meristem}{形成层}形成的不定芽；每个芽都是一个完整的\omterm{def:b2:plant-meristems:meristem}{分生组织}，能够
重建枝条系统，而动物的头部是一个单一的
不可替换的\omterm{def:b2:vertebrate-development:induction}{组织者}，没有备用副本。
\textbf{24.} 到达芽处的\omterm{prop:b2:plant-meristems:apical}{细胞分裂素}增多，促进芽长出
——给修剪过的植株浇水，使其分枝更多。
\textbf{25.} 每个顶端 50 片叶，每棵树 \num{10000} 片；节间伸长速率
白天为 \qty{0.01}{h^{-1}}，夜晚为 \qty{0.02}{h^{-1}}；全年
\omterm{def:b2:plant-meristems:wood}{木材} \qty{0.025}{m^3}，干重 \qty{12.6}{kg}，其中碳 \qty{6.3}{kg}。
\end{solution}
